\documentclass[11pt]{book}
\usepackage[utf8]{inputenc}
\usepackage{geometry}
\geometry{a5paper, margin=1.2in}
\usepackage{ebgaramond}
\usepackage{titlesec}
\usepackage{epigraph}
\usepackage{xcolor}
\usepackage{setspace}
\onehalfspacing
\setlength{\epigraphwidth}{0.7\textwidth}
\definecolor{ashgray}{gray}{0.3}
\newcommand{\ledger}[1]{\textsc{#1}}
\newcommand{\ritetag}[1]{\textsf{[#1]}}
\newcommand{\clock}[1]{\ledger{#1 Clock}}

\titleformat{\chapter}{\normalfont\huge\itshape}{\thechapter}{20pt}{}
\titleformat{\section}{\normalfont\Large\scshape}{\thesection}{15pt}{}

\begin{document}

\frontmatter
\title{\Huge\itshape Bitter Root\\\large Tales of Unseen Labor, Female Rage, and the Cords That Bind}
\author{Dusana of the Raven Road}
\date{Under the lamp, bread and salt——the ink is ash and well-water that remembers.}
\maketitle

\chapter*{Preface: The Root That Feeds, the Root That Binds}

\epigraph{The women who wash the dead do not sing. They count. One for the breath that left, two for the name that stays, three for the grief they will carry home. No one thanks them. The soup is free. The price is never asked.}{——Old Mila, hedge- keeper of the Moss Road}

A ledger of stories is a ledger of unpaid hours. The washing, the nursing, the midwifery, the stirring of pots that never empty——none of it appears in the sagas. Heroes do not mend fences. Queens do not change linens. And yet the world turns because someone, somewhere, is doing the work that leaves no mark.

This book is not for heroes. It is for the women who stay in the kitchen while the baronet feasts, who walk into the fog to find a lost child, who pour the eighth cup and leave the ninth empty because they know the Hollow drinks last. It is for the rage that does not scream——the rage that braids itself into patience and calls itself \textit{care}.

The first story is a flight. A village burns, a girl runs, and a coven of hedge-witches, Daughter renegades, and a Tulkani storyteller smuggle her across three borders. The second story is a siege by other means: a Hearth Witch and a Daughter of the Chain, locked in a cold war over a baronet's household, who slowly recognize each other's teeth and decide——grudgingly, frostily——to work together.

I have written them in the old way: by a borrowed lamp, on paper that had been a bandage. The mechanics are here, but they are not the point. The point is the weight of the work. The point is the thread.

— Dusana of the Raven Road, at a crossroads where the well still sings

\mainmatter

\chapter{The Smell of Thyme}
\epigraph{They burned the orchard first. Thyme grows back. Girls don't.}{——Elara, age twelve, before she forgot her name}

\section*{I. The Root Cellar}

The smell of thyme is the last thing Elara remembers clearly. Her mother had planted it along the eastern wall, where the morning sun hit the stones just so. \textit{Thyme for courage}, her mother used to say, pinching a leaf between her fingers. \textit{Thyme for memory.}

Then the torches came over the hill.

Elara is twelve, small for her age, with a mage's knot behind her left ear that she has never learned to name. She can hear the last words of dying plants. It is not a useful gift——until it is. When the wheat screams, you know the fire is close.

The Chain-Lanterns of Thepyrgos move with the precision of men who have done this before. Five of them, hooded in iron mesh, their leader carrying a writ stamped with a patriarchal seal. They do not knock. They kick in the door of the house where Elara's mother is still cutting bread.

\textit{Where is the girl?}

Her mother says nothing. She is good at nothing. It is the first gift a village witch gives her daughter: the ability to become furniture.

Elara is already in the root cellar, the door concealed behind a rack of drying herbs. She counts her breaths——eight, never nine, because nine is for the Hollow. Through the cracks in the floorboards, she sees the glint of a Lantern's bell. They ring it once before accusation, three times before burning. She has heard that bell in her nightmares for years. Now it is real.

\section*{II. Old Mira's Last Lesson}

The village witch, Old Mira, is not surprised. She has been expecting the Chain-Lanterns since the autumn, when a neighbor's son saw Elara whispering to a frost-killed apple tree and heard the tree whisper back. Mira has already sent word to the Breath-less network——a cracked clay bowl left on the north doorstep, a red thread tied around the well rope.

She finds Elara in the root cellar as the first thatch catches fire.

\epigraph{Draw a line. Any spirit will respect a line if you draw it with witness.}{——Old Mira, her last words}

Mira gives Elara three things: a piece of chalk, a knot of red thread, and a name that is not her own. \textit{You are not Elara anymore. You are Thyme. And you will run.}

The Chain-Lanterns' procedure is relentless. They establish a \ledger{Lantern Line} around the village——women's rites cannot cross without test. They demand that every household present its patriarch to swear loyalty. Old Mira, who has no patriarch, who has been the village's memory for forty years, is dragged into the square.

They do not burn her immediately. They question her first, because the Inquisition loves procedure. \textit{Where is the girl? Name her true name.}

Mira looks at the smoke rising from the orchard. The thyme is burning. She smiles.

\textit{She is already gone. And you will not remember her face.}

It is not a spell. It is a curse, and she has been weaving it for years. Every bowl of soup she gave the village, every whispered comfort at a deathbed, every child she delivered——all of it was thread. Now she pulls.

The Chain-Lanterns' leader blinks. He knows Elara's mother's face, but the daughter's? It slips from his memory like water from a cracked bowl. He writes \textit{girl, approximately twelve, dark hair} in his ledger, and even as he writes, he knows it is not enough.

\section*{III. The Covens' Road}

Elara——she will learn to answer to Thyme, but not yet——runs through the burning orchard. The thyme smoke is thick and green. It covers her scent. Behind her, the village screams. Ahead, the dark.

The first safe house is a charcoal-burner's hut three miles into the woods. The woman there is not a witch. She is a widow who lost her daughter to a fever and now owes the Breath-less a debt she will never fully repay. She gives Elara bread and salt, and a new pair of boots that are too large.

\textit{The Hearth-Menders will meet you at the ford.}

The Hearth-Menders are not a coven in the formal sense. They have no patron, no thiasos, no written code. They have clay bowls and lullabies sung in harmonics, and they understand one thing: \textbf{injury obligates repair}. A village that burns owes the survivors a future.

\subsection*{Mechanics: Hearth-Mender Hospitality}
When the Hearth-Menders shelter a fugitive, they create a \clock{Task-Bond} of 4 segments. The fugitive must perform a small act of care——mending a net, telling a true story, watching a child for an hour——before receiving the next safe passage. This is not cruelty; it is the law of mutual aid. \textit{Share your hunger; name your fear; accept a task before a gift.}

Elara does not understand why she must braid the widow's hair before the widow will give her a cloak. But she does it. Her hands shake. The widow's hair smells of woodsmoke. When she finishes, the widow says, \textit{Now you are not just a girl running. You are someone who gave something. That is harder to forget.}

\section*{IV. The Daughter's Ferry}

The black-water river is the border. On the far side, the Chain-Lanterns' writ means nothing. But the river itself is a threshold, and thresholds demand tolls.

The ferry is a flat-bottomed boat, its pilot a woman in grey robes with a brass hook for a hand and a child standing beside her——a Hollowed, its eyes like spilled milk. A Daughter of the Chain. A renegade.

Elara hesitates. Everyone knows the Daughters. The soup is free. The price never is.

\epigraph{Give me a memory. Not a pretty one. The one you would least like to lose.}{——Sister Bryn of the Unbroken Chain (renegade)}

The Daughter's name is Bryn——Sister Bryn, though she no longer wears the cord of the coven. She left the High Mother's service after she was ordered to harvest a dying woman's last breath and the woman had been kind to her. There is a clock ticking in Bryn's flesh, a \ledger{Redemption} clock that she does not know how to fill.

She rows in silence. The Hollowed watches Elara with its milk-white eyes.

\textit{The smell of thyme,} Elara says. \textit{That's the memory. My mother's garden.}

Bryn nods. She does not write it down. She does not need to. The memory is now a knot in her own Cord, a debt she will carry until she finds a way to pay it forward.

On the far bank, Elara stumbles ashore. Bryn does not say goodbye. She turns the boat and rows back into the fog, the Hollowed's hand on her sleeve.

\subsection*{Mechanics: Renegade Daughter's Cord}
A Daughter who breaks from the Chain keeps her Hollowed but no longer owes the High Mother. Instead, she tracks a \clock{Redemption} of 6 segments. Each time she gives shelter or aid without demanding a price, she marks a segment. When the clock fills, she may choose to destroy her Hollowed, rejoin the coven, or transform her Cord into a new kind of bond——one built on mercy rather than debt.

\section*{V. The Silk-Weaver's Attic}

The last safe house is a Lethai-ar silk-weaver's cottage in the Crimson Valley. The weaver, a woman called \textit{Mask-Quill}, does not speak to Elara for three days. She simply sets a bowl of soup on the table each morning and takes the empty bowl away each evening.

On the fourth day, she asks: \textit{What is your name?}

\textit{Elara.}

\textit{No. That name is ash. Choose another.}

Elara looks at the bowl. The soup is thick with thyme.

\textit{Thyme,} she says. \textit{My name is Thyme.}

Mask-Quill nods. She takes a piece of red silk and ties it around Thyme's wrist. \textit{This is a witness cord. It will remember the truth of you, even when you forget.}

She does not explain what the truth is. Thyme does not ask. She is twelve. She has watched her village burn. She has crossed a black river on a Daughter's boat. She has learned that the world runs on debts she did not incur and must repay anyway.

She will not forget the smell of thyme.

\section*{VI. The Bitter Root}

\epigraph{The root that feeds is the root that binds. The hand that rocks the cradle is the hand that sharpens the knife.}{——from the \textit{Book of the Unnamed}, forbidden in Thepyrgos}

Elara——Thyme will grow up in the Crimson Valley. She will learn to weave silk and to tie witness cords. She will never see her mother again. She will never know that Old Mira's last curse worked: the Chain-Lanterns' report described a girl with brown hair and no name, and the Inquisition closed the case six months later.

But she will keep the red thread on her wrist. And when she is old enough, she will join the Hearth-Menders. Not out of revenge. Out of \textit{care}. Because someone has to stir the pot.

The soup is free. The second bowl will cost you a story.

\subsection*{Mechanics: Bitter Root Condition}
A character who has survived state violence and been smuggled to safety may acquire the \ledger{Bitter Root} condition. It grants +1 die to acts of mutual aid (feeding, sheltering, healing others) but –1 die to acts of self-advocacy (asking for pay, demanding recognition, refusing a burden). The condition can be \textit{pruned} by an act of witnessed rage——a public accusation, a broken tool, a screamed truth. Pruning removes the penalty and converts the bonus to a permanent \ledger{Coven Bond} String.

\newpage

\chapter{The Second Hearth}
\epigraph{She stirred the pot. She watched the pot. They never saw her add the salt.}{——Morwen of Thornwood, on the occasion of no thanks whatsoever}

\section*{I. The New Baronet}

Thornwood is a village of sixty souls, a mill, a ford, and a Hearth Witch named Morwen. She has served the community for thirty years——mending bones, settling disputes, keeping the Hollow from the door. She asks nothing but a bowl of soup and a place by the fire. The villagers do not thank her. That is fine. Thanking is not why she does it.

Then the crown appoints a new Knight-Baronet: Sir Aldric of the Eastern Marches, a man with a hunting hawk, a young wife, a sickly infant, and no understanding of hedges. He arrives in a carriage that gets stuck in the ford. His escort tramples the miller's bean field. He does not apologize.

Morwen watches from the kitchen door of her cottage. She already dislikes him.

But it is the wetnurse who makes her blood run cold.

The woman is tall, grey-haired, efficient. She moves like a woman who has attended a thousand births and a thousand deaths. She wears no obvious mark of the Chain, but Morwen has been a threshold worker for three decades. She knows the signs: the way the woman never blinks when a baby cries, the way she always pours tea \textit{first for the master}, the way she has her own set of scales for measuring milk.

Her name is Sister Bryn.

\epigraph{You're the one who left the Chain. I heard stories.}{——Morwen, the first night}

\epigraph{Stories are debts. I'm not here to pay yours.}{——Sister Bryn, the first night}

\section*{II. The Cold War}

The baronet's household is a small court: Sir Aldric, his wife Lady Margot, the infant Eleanor, a cook, a groom, and Bryn. Morwen is not part of the household. She is the village witch, which means she is expected to help when things go wrong and to stay out of sight when things go right.

The cold war begins with small things.

Bryn reorganizes Morwen's herb stores \textit{for efficiency}. Morwen misplaces Bryn's witness cords \textit{to air them out}. Bryn salts the well water (just a pinch, not enough to taste, enough to make Morwen's divinations cloudy). Morwen ties a red thread around the kitchen door——a \ledger{Warm Hand} working that makes the room feel too hot for anyone wearing a Daughter's grey.

Neither woman speaks of it. The baronet notices nothing. His wife notices everything but says nothing.

\subsection*{Mechanics: Household Disposition}
A dwelling under conflicting influences gains a \ledger{Contested} disposition. Threshold workings inside suffer +1 DV. Social rolls involving hospitality are Dominant if you side with one faction, Desperate if you try to please both. The disposition changes only when one witch withdraws or a crisis forces cooperation.

\section*{III. The Infant's Fever}

Lady Margot's baby, Eleanor, is a sickly thing. She coughs, she does not gain weight, she sleeps too much. Morwen has suspected for weeks that the child has a weak heart——something no poultice can fix, only careful nursing.

One night, the fever spikes. Lady Margot comes pounding on Morwen's door at midnight.

\textit{Please. She's burning up.}

Morwen goes. She brings feverfew, a silver spoon, and her own mother's lullaby. She finds Bryn already there, holding the infant, her face as still as a mask.

The two women work side by side. Morwen cools the child's forehead with a damp cloth while Bryn measures out drops of a grey tincture. Neither speaks. The only sounds are the infant's ragged breath and the soft chime of the baronet's hunting clock in the next room.

At dawn, the fever breaks.

\epigraph{Tincture of willow bark and a drop of your own blood. You're funding the Cord with her survival.}{——Morwen, accusatory}

\epigraph{And your lullaby is a prayer to Inaea. You're binding her to gratitude she'll never be able to repay. Same coin, different face.}{——Bryn, equally accusatory}

They stare at each other across the crib. The infant sleeps, her cheeks pink for the first time in weeks.

\textit{I didn't do it for the debt,} Morwen says.

\textit{Neither did I. But the debt accrues anyway.}

\section*{IV. The Missing Child}

The winter that follows is harsh. The ford freezes. The mill's waterwheel cracks. Three villagers come down with a lung-sickness that Morwen recognizes as the \ledger{Spreading}——a Daughter's curse, she is sure of it. But Bryn tends the sick alongside her, and the sickness fades without a single death.

Then a village child——a girl named Pella, seven years old, curious and quick——wanders into the woods at twilight and does not return.

The baronet organizes a hunting party. Morwen knows the forest; Bryn can read the Hollow's attention. They go together, leaving the men to thrash about with torches.

\subsection*{Mechanics: Cooperative Threshold Working}
Two witches working toward the same goal can combine their Threshold talents. One rolls \ledger{Lore} (to track the Hollow's signs), the other rolls \ledger{Survival} (to navigate the forest). On a success, they reduce a \clock{Search} clock by 2 segments. On a failure, the Hollow notices them——gain 1 \ledger{Hollow Attention} tick.

They find Pella at midnight, sitting by a standing stone, unharmed. The Hollow was \textit{just looking}, the girl says. \textit{It was lonely.}

Morwen takes the child's left hand. Bryn takes the right. They walk out of the woods without speaking. Behind them, the standing stone hums once, then falls silent.

\section*{V. The Winter Kitchen}

The baronet throws a feast at midwinter——roasted swan, spiced wine, a fire so large it eats a month's worth of wood. He does not invite Morwen. He does not thank Bryn. He does not notice that the food on the table came from the village's stores, or that the wine was watered, or that the servants are wearing the same coats they wore last year.

In the kitchen, Morwen stirs a pot of soup for the villagers who were not invited. Bryn portions the bread.

\textit{You're not so useless,} Morwen says.

\textit{You're not so sentimental,} Bryn replies.

They drink a cup of small beer. They do not toast. They do not smile. But they do not glare.

\epigraph{The red thread on the kitchen door. Did you tie that?}{——Bryn}

\epigraph{I thought you did.}{——Morwen}

Neither woman unties it.

\section*{VI. The Thaw}

Spring comes. The ford opens. The baronet receives a letter recalling him to court——something about a border dispute, nothing that concerns him. He leaves his wife and child behind. He does not say goodbye to the village.

Lady Margot, to everyone's surprise, asks Morwen to teach her about herbs. \textit{I want to be useful,} she says. Morwen agrees, on the condition that Bryn teaches her about the Hollow in return.

The three women eat together once a week: Morwen, Bryn, and Lady Margot. They do not become friends. Their conversation is the conversation of people who have seen too much to pretend. But they build a small, fragile thing: a community that is not based on blood or marriage or feudal obligation, but on the simple recognition that \textit{someone has to do the work}.

\subsection*{Mechanics: Unseen Labor Clock}
A community that is well tended——by witches, midwives, hedgers, and nurses——maintains a \ledger{Community Health} clock of 6 segments. Each week of care ticks one segment. If the clock is full at the end of winter, the community survives any hardship without loss of life. If it empties, someone starves, or freezes, or sickens beyond help.

Morwen and Bryn, working together, keep the clock full. They never tell the baronet. He never asks.

\section*{VII. The Final Scene}

One evening, as the last snow melts, Morwen finds a single red thread tied between a chair leg and the pot handle. She knows she did not tie it. She knows Bryn did not tie it.

She leaves it there.

When Bryn comes to fetch the soup pot, she sees the thread. She does not mention it.

But she pours two cups of tea instead of one.

\epigraph{The soup is free,}{——Morwen, pushing a cup across the table}

\epigraph{The price is never,}{——Bryn, taking it}

They drink in silence. The pot simmers. The thread holds.

\chapter*{Colophon}

This book was written on paper that had been a grain sack, in ink cut with ash from a Chain-Lantern's torch. The binding is red thread, tied with eight knots.

The ninth knot is not cut. The work is not done. The women who wash the dead are still counting.

— Dusana of the Raven Road, at a crossroads where three covens meet

\appendix
\chapter{Mechanics Summary}

\section*{Bitter Root Condition}
\ledger{Bitter Root} (persistent): +1 die to mutual aid actions, –1 die to self-advocacy. Pruned by an act of witnessed rage.

\section*{Unseen Labor Clock}
6 segments. Ticks with weekly care. Full = community survives hardship. Empty = loss.

\section*{Renegade Daughter's Cord}
\clock{Redemption} (6 segments). Ticks when a Daughter gives aid without demanding price. Full grants choice: destroy Hollowed, rejoin Chain, or transform Cord.

\section*{Cooperative Threshold Working}
Two witches combine \ledger{Lore} and \ledger{Survival} to reduce a \clock{Search} clock by 2 on success. On failure, gain \ledger{Hollow Attention}.

\section*{Household Disposition}
\ledger{Contested}: threshold workings +1 DV; social rolls Dominant if siding with one faction, Desperate if trying to please both.

\end{document}
